\documentclass[12pt, a4paper]{article}
\usepackage[utf8]{inputenc}
\usepackage[T2A]{fontenc}
\usepackage[russian]{babel}

\usepackage{amsmath,amsfonts,amssymb,amsthm,mathtools}
\usepackage{amssymb}

\begin{document}

\textbf{Здесь записана формула:}

\begin{equation}
\begin{split}
f = x + x \cdot x
\end{split}
\end{equation}

\textbf{Оптимизация формулы:}
\begin{align*}
f = x + x \cdot x\end{align*}

\textbf{Дифференцирование формулы:}
\begin{align*}
 \frac {\partial f} {\partial x}  = 1 + 1 \cdot x + x \cdot 1\end{align*}

\textbf{И снова оптимизация формулы:}
\begin{align*}
 \frac {\partial f} {\partial x}  = 1 + x + x\end{align*}

\textbf{Дифференцирование формулы:}
\begin{align*}
 \frac {\partial^2 f} {\partial x \cdot \partial x}  = 0 + 1 + 1\end{align*}

\textbf{И снова оптимизация формулы:}
\begin{align*}
 \frac {\partial^2 f} {\partial x \cdot \partial x}  = 2\end{align*}

\textbf{Дальнейшие преобразования, оставим читателю в качестве самостоятельного упражнения}.

Approved by "Кафедра вышмата"
\end{document}
